\section{Training a policy in Wellocate}
\label{app:b11}

Trained on Wellocate's reward alone, a two-parameter rule already settles on a setting that is best or tied for
best on generated landscapes (Appendix~\ref{app:sim}). This appendix trains a full policy.

\paragraph{Design.} We trained a small policy with proximal policy optimisation (PPO) through Wellocate's released
interface. It sees what a language model agent sees: the readings, the outcomes and the remaining budget, and never a
true value or the set of viable candidates. At each step it scores a set of candidates. The set holds re-tests of
what it has ordered, every single-letter change to the start and to its eight best candidates, and one random untried
candidate. A second network picks the nominee and a third estimates value. Its inputs are scaled to the campaign's own
readings, so the same networks run on any landscape. Run A trains on generated grid landscapes only. Run B also
trains on scrambled GB1 and is tested on PhoQ. Each run has three seeds and keeps the checkpoint with the best score
on its validation split. The design, the test and the success criteria were registered before training. The policy had to beat both its own untrained
initialisation and random ordering, each with a 95\% interval above zero and every seed positive. It had to do so on
the grid test split to show the task is learnable, and on the scrambled GB1 and PhoQ test splits to show transfer.

\paragraph{The registered result.} Run A meets both criteria (Table~\ref{tab:b11reg}). Run B beats its initialisation,
and its interval against random ordering includes zero, so it fails the transfer criterion. A rerun decided after the
test, with an exploration bonus five times the registered one, meets it on PhoQ. It beats random ordering by
\nPolEntRandPq{} \nPolEntRandPqci{} with every seed positive (Table~\ref{tab:b11phase2}). The registered intervals
resample episodes only. Table~\ref{tab:b11reg} gives every interval from a two-level bootstrap that resamples
landscapes and then episodes; this check is post-hoc. Run A's criteria all hold under it, and also under an interval
across the three seeds. Under the two-level bootstrap, two comparisons with the guessing rule include zero, on the grid and on PhoQ, so we make no
claim that the policy beats the guessing rule. Five more seeds of Run A, trained and tested the same way, replicate
this: every seed beats its initialisation and random ordering on all three test splits (Table~\ref{tab:b11phase2}).

\paragraph{Robustness.} Everything from here on was decided after the registered test. Run A plays every condition without retraining, with
the networks of the registered test. Its lead over random ordering holds when the noise is halved or doubled
(Table~\ref{tab:b11robust}). A per-round batch effect, a noise floor fifty times the smallest spread of Appendix~\ref{app:env} and failed wells leave the
scripted policies' best number of re-tests unchanged. Under each of them, three or four re-tests a round stay below
that best (Table~\ref{tab:b11stress}). Spikes on the grid are the one exception. When 2\% of readings are three times their
value, the informed rule with three re-tests a round becomes the best scripted rule. Its pick minus that of its best
setting with fewer re-tests is \nPolSpikeHeavy{} \nPolSpikeHeavyci. Run A does not adapt to the spikes. It keeps
re-testing about \nPolShare\% of its slots, and its lead over random ordering falls to \nPolRandSpike{}
\nPolRandSpikeci{} (Table~\ref{tab:b11robust}). On GB1 and PhoQ under their real names it behaves as on the scrambled
test splits (Table~\ref{tab:b11robust}). Its lead shrinks on rougher grids and holds at every budget we tried
(Table~\ref{tab:b11robust}). At 4 rounds of 24 it is level with the best scripted rule, with two of its inputs beyond the range it saw in training.

\paragraph{The language models' own campaigns.} We replayed Run A on each model's own campaigns: the same
landscape, campaign seeds, noise and budget (Table~\ref{tab:b11llm}). We first checked that each of these landscapes
is outside the training and validation splits, by name, by built content and, for the grid, by its sorted fitness values. On
the grid Run A beats every model. On scrambled GB1 it is below GPT-6.1 Sol, and on GB1 under real names it is below
Opus 4.7.

\paragraph{What the policies learned.} We replayed every test episode and read each order back from the
environment's record (Table~\ref{tab:b11behaviour}). Run A builds from its leaders and rarely re-tests them. On the
grid, \nPolParentTopThree\% of its new candidates are one letter from one of its three highest readings. Re-tests of
those three take \nPolLeadRetest\% of its orders. Most of its re-tests retry an order that failed, \nPolRetry\%, or
repeat a candidate first read earlier the same round, \nPolSameRound\%. The seeds differ: the third seed re-tests candidates that already have a reading more often. The registered Run B re-tests from the first round, and its share climbs to about four slots in five by the
last. Most of its re-tests go to candidates it has read. On GB1 Run A's nominee is often not its highest reading, and
that costs it much of its gap to the best scripted rule there.

% Generated by scripts/make_table_b11.py from results/b11_*.json. Do not edit by hand;
% scripts/verify_paper_numbers.py regenerates it and fails on any difference.

\begin{table}[htbp]
\centering
\small
\setlength{\tabcolsep}{1.5pt}
\caption{The registered test of the trained policy, 1{,}500 episodes per split, as a mean over three seeds. Run A trained on generated landscapes only; Run B on generated landscapes and scrambled GB1, tested on PhoQ. Differences are in final pick, with 95\% intervals from a two-level bootstrap that resamples landscapes and then episodes (post-hoc; the registered intervals resample episodes only).}
\label{tab:b11reg}
\begin{tabular}{lrrrr}
\toprule
 & Run A, grid & Run A, GB1 & Run A, PhoQ & Run B, PhoQ \\
\midrule
trained pick & 7.91 & 4.79 & 2.15 & 1.53 \\
re-testing share (\%) & 9.3 & 1.8 & 2.3 & 68.3 \\
minus initialisation & $+$0.86 [0.83, 0.89] & $+$3.32 [3.24, 3.40] & $+$1.49 [1.44, 1.54] & $+$0.87 [0.84, 0.91] \\
minus random & $+$0.40 [0.36, 0.44] & $+$1.43 [1.34, 1.53] & $+$0.65 [0.58, 0.73] & $+$0.03 [$-$0.03, 0.09] \\
minus guessing rule & $-$0.03 [$-$0.06, 0.01] & $-$1.32 [$-$1.45, $-$1.20] & $+$0.07 [$-$0.01, 0.15] & $-$0.55 [$-$0.62, $-$0.48] \\
minus informed rule & $-$0.05 [$-$0.10, 0.00] & $-$1.16 [$-$1.26, $-$1.07] & $-$0.52 [$-$0.63, $-$0.41] & $-$1.14 [$-$1.25, $-$1.03] \\
minus model-guided & $-$0.06 [$-$0.11, $-$0.01] & $-$0.56 [$-$0.67, $-$0.45] & $+$0.27 [0.21, 0.33] & $-$0.35 [$-$0.39, $-$0.31] \\
\bottomrule
\end{tabular}
\end{table}

\begin{table}[htbp]
\centering
\scriptsize
\setlength{\tabcolsep}{3pt}
\caption{The trained policy (Run A, without retraining) under every change after the registered test, 1{,}500 episodes per row, at assay noise unless the row changes it. Rows that would repeat Table~\ref{tab:b11reg} are left out. The last two columns are Run A's pick minus random ordering and minus the best scripted rule at its best number of re-tests, on the same episodes. Stress: each round's readings share one multiplicative factor with this standard deviation (batch), an additive noise floor of 2\% of the landscape maximum (floor), or 2\% of readings replaced by a near-zero value (failed wells) or by three times the reading (spikes). Real names: GB1 and PhoQ under their amino-acid names. At 4 rounds of 24 the scripted rules' re-tests scale with the slots per round, and two of Run A's input features exceed the range seen in training. Intervals as in Table~\ref{tab:b11reg}; under real names, one landscape each, they resample episodes only.}
\label{tab:b11robust}
\begin{tabular}{llrrrr}
\toprule
change & condition & pick & share (\%) & minus random & minus best scripted rule \\
\midrule
noise, grid & noise halved (3.65\%) & 8.11 & 9.2 & $+$0.45 [0.40, 0.49] & $-$0.04 [$-$0.08, 0.00] \\
 & noise doubled (14.6\%) & 7.66 & 9.0 & $+$0.35 [0.31, 0.40] & $-$0.09 [$-$0.14, $-$0.05] \\
\midrule
stress, grid & batch 5\% & 7.84 & 9.3 & $+$0.41 [0.36, 0.46] & $-$0.13 [$-$0.18, $-$0.07] \\
 & batch 10\% & 7.73 & 9.2 & $+$0.40 [0.35, 0.45] & $-$0.13 [$-$0.19, $-$0.08] \\
 & floor & 7.89 & 9.3 & $+$0.40 [0.35, 0.44] & $-$0.11 [$-$0.16, $-$0.06] \\
 & failed wells & 7.90 & 9.5 & $+$0.39 [0.35, 0.43] & $-$0.10 [$-$0.14, $-$0.06] \\
 & spikes & 7.46 & 9.1 & $+$0.09 [0.05, 0.13] & $-$0.21 [$-$0.24, $-$0.17] \\
\midrule
stress, GB1 & batch 5\% & 4.78 & 1.9 & $+$1.43 [1.33, 1.53] & $-$1.23 [$-$1.35, $-$1.11] \\
 & batch 10\% & 4.73 & 1.8 & $+$1.40 [1.28, 1.51] & $-$1.22 [$-$1.35, $-$1.09] \\
 & floor & 4.82 & 1.9 & $+$1.47 [1.37, 1.58] & $-$1.26 [$-$1.38, $-$1.14] \\
 & failed wells & 4.75 & 1.8 & $+$1.42 [1.32, 1.53] & $-$1.30 [$-$1.44, $-$1.16] \\
 & spikes & 4.44 & 1.8 & $+$1.17 [1.06, 1.27] & $-$1.05 [$-$1.19, $-$0.92] \\
\midrule
real names & GB1 & 4.72 & 1.9 & $+$1.31 [1.24, 1.39] & $-$1.37 [$-$1.46, $-$1.29] \\
 & PhoQ & 2.12 & 2.3 & $+$0.58 [0.53, 0.63] & $-$0.47 [$-$0.54, $-$0.39] \\
\midrule
held-out grids & $K=2$ & 7.70 & 8.7 & $+$0.18 [0.12, 0.23] & $-$0.27 [$-$0.31, $-$0.23] \\
 & $K=3$ & 7.55 & 8.3 & $+$0.04 [$-$0.03, 0.11] & $-$0.31 [$-$0.37, $-$0.25] \\
\midrule
budgets, grid & 10 rounds of 4 & 7.83 & 9.1 & $+$0.42 [0.36, 0.47] & $-$0.06 [$-$0.12, $-$0.01] \\
 & 20 rounds of 4 & 7.97 & 9.5 & $+$0.40 [0.36, 0.44] & $-$0.11 [$-$0.15, $-$0.07] \\
 & 4 rounds of 24 & 8.01 & 11.6 & $+$0.37 [0.32, 0.43] & $-$0.04 [$-$0.09, 0.00] \\
\bottomrule
\end{tabular}
\end{table}

\begin{table}[htbp]
\centering
\scriptsize
\setlength{\tabcolsep}{3pt}
\caption{What the trained policies do on the test splits, as a mean over three seeds. Share: re-tests as a share of the budget. The next three columns split the re-tests: retries of an order that failed, repeats of a candidate first read earlier the same round, and re-tests of one of the three candidates with the highest mean reading. Parent: new candidates one letter from one of those three. Nominee: how often it is the registered selection rule's nominee, and its pick minus that nominee's.}
\label{tab:b11behaviour}
\begin{tabular}{llrrrrrrr}
\toprule
run & split & share (\%) & retries (\%) & same round (\%) & top three (\%) & parent (\%) & nominee (\%) & pick cost \\
\midrule
Run A & grid & 9.3 & 62 & 15 & 5 & 61 & 97 & $+$0.00 [0.00, 0.00] \\
Run A & GB1 & 1.8 & 58 & 29 & 4 & 63 & 51 & $-$0.71 [$-$0.75, $-$0.67] \\
Run A & PhoQ & 2.3 & 61 & 23 & 2 & 64 & 89 & $-$0.05 [$-$0.06, $-$0.04] \\
\midrule
Run B & grid & 72.0 & 3 & 5 & 58 & 43 & 71 & $-$0.03 [$-$0.04, $-$0.02] \\
Run B & GB1 & 69.5 & 3 & 5 & 42 & 35 & 93 & $-$0.06 [$-$0.07, $-$0.05] \\
Run B & PhoQ & 68.3 & 3 & 3 & 39 & 31 & 93 & $-$0.02 [$-$0.02, $-$0.01] \\
\bottomrule
\end{tabular}
\end{table}

\begin{table}[htbp]
\centering
\scriptsize
\setlength{\tabcolsep}{2pt}
\caption{Five more seeds of Run A (three to seven), trained and tested as registered, and Run B retrained with an entropy bonus five times the registered one, tested with the registered code on all three test splits. Run B trains on the grid and GB1, so only PhoQ tests its transfer. Differences in final pick against random ordering, with intervals as in Table~\ref{tab:b11reg}. Entropies are at the end of training: means over the last updates, averaged over seeds.}
\label{tab:b11phase2}
\begin{tabular}{lrrr}
\toprule
 & Run A, seeds 3 to 7 & Run B, registered & Run B, larger bonus \\
\midrule
grid: minus random & $+$0.42 [0.37, 0.46] & -- & $+$0.32 [0.28, 0.37] \\
grid: re-testing share (\%) & 10.3 & -- & 3.2 \\
grid: seeds above random & 5 of 5 & -- & 3 of 3 \\
GB1: minus random & $+$1.63 [1.53, 1.72] & -- & $+$1.66 [1.56, 1.76] \\
GB1: re-testing share (\%) & 1.7 & -- & 2.7 \\
GB1: seeds above random & 5 of 5 & -- & 3 of 3 \\
PhoQ: minus random & $+$0.66 [0.59, 0.72] & $+$0.03 [$-$0.03, 0.09] & $+$0.50 [0.43, 0.56] \\
PhoQ: re-testing share (\%) & 1.9 & 68.3 & 3.1 \\
PhoQ: seeds above random & 5 of 5 & 2 of 3 & 3 of 3 \\
order-head entropy, nats & 5.31 & 1.06 & 5.88 \\
nominate-head entropy, nats & 0.04 & 0.03 & 0.34 \\
\bottomrule
\end{tabular}
\end{table}

\begin{table}[htbp]
\centering
\scriptsize
\setlength{\tabcolsep}{3pt}
\caption{The trained policy (Run A, without retraining) on each language model's own campaigns, the same landscape, campaign seeds, noise and budget. Pick, with re-testing share; the last column is Run A's pick minus the model's, paired by campaign, with a 95\% bootstrap interval over campaigns. Each of these landscapes is outside the training and validation splits.}
\label{tab:b11llm}
\begin{tabular}{llrrrr}
\toprule
landscape & model & campaigns & model & Run A & Run A minus model \\
\midrule
grid & Opus 4.7 & 60 & 7.62 (75\%) & 7.87 (10\%) & $+$0.25 [0.14, 0.35] \\
grid & Gemini 2.5 Flash & 120 & 7.70 (94\%) & 7.86 (9\%) & $+$0.16 [0.11, 0.21] \\
grid & DeepSeek V3.1 & 120 & 7.76 (91\%) & 7.86 (9\%) & $+$0.10 [0.04, 0.15] \\
grid & GPT-6.1 Sol & 60 & 7.67 (44\%) & 7.87 (10\%) & $+$0.19 [0.07, 0.31] \\
grid & Gemini 3.1 Pro & 60 & 7.69 (14\%) & 7.87 (10\%) & $+$0.18 [0.07, 0.29] \\
\midrule
GB1, scrambled & Opus 4.7 & 60 & 4.78 (63\%) & 4.79 (2\%) & $+$0.02 [$-$0.37, 0.42] \\
GB1, scrambled & Gemini 2.5 Flash & 60 & 1.12 (92\%) & 4.79 (2\%) & $+$3.68 [3.42, 3.92] \\
GB1, scrambled & DeepSeek V3.1 & 60 & 0.84 (96\%) & 4.79 (2\%) & $+$3.95 [3.67, 4.22] \\
GB1, scrambled & GPT-6.1 Sol & 30 & 6.36 (19\%) & 4.67 (2\%) & $-$1.69 [$-$2.09, $-$1.35] \\
GB1, scrambled & Gemini 3.1 Pro & 30 & 4.90 (11\%) & 4.67 (2\%) & $-$0.24 [$-$0.82, 0.37] \\
\midrule
GB1, real names & Opus 4.7 & 60 & 7.59 (51\%) & 4.75 (2\%) & $-$2.84 [$-$3.21, $-$2.46] \\
GB1, real names & Gemini 2.5 Flash & 60 & 2.41 (45\%) & 4.75 (2\%) & $+$2.34 [1.93, 2.74] \\
GB1, real names & DeepSeek V3.1 & 60 & 1.27 (92\%) & 4.75 (2\%) & $+$3.48 [3.11, 3.84] \\
\midrule
PhoQ, scrambled & Opus 4.7 & 30 & 1.89 (63\%) & 2.20 (2\%) & $+$0.31 [0.06, 0.58] \\
PhoQ, scrambled & Gemini 2.5 Flash & 30 & 2.13 (64\%) & 2.20 (2\%) & $+$0.07 [$-$0.16, 0.33] \\
PhoQ, scrambled & DeepSeek V3.1 & 30 & 0.83 (90\%) & 2.20 (2\%) & $+$1.37 [1.17, 1.60] \\
PhoQ, scrambled & GPT-6.1 Sol & 30 & 2.37 (24\%) & 2.20 (2\%) & $-$0.17 [$-$0.44, 0.12] \\
PhoQ, scrambled & Gemini 3.1 Pro & 30 & 1.72 (15\%) & 2.20 (2\%) & $+$0.48 [0.24, 0.72] \\
\bottomrule
\end{tabular}
\end{table}

\begin{table}[htbp]
\centering
\scriptsize
\setlength{\tabcolsep}{3pt}
\caption{The guessing and informed rules under the stress conditions of Table~\ref{tab:b11robust} that keep their best number of re-tests, 1{,}500 episodes per landscape and condition, at assay noise. Without stress their best is 1 re-test a round on the grid and 0 on GB1. Best k is the best number of re-tests a round, with its re-testing share; three minus best is the pick of three re-tests a round minus that of the best, with intervals as in Table~\ref{tab:b11reg}. Four re-tests a round are lower still in every row.}
\label{tab:b11stress}
\begin{tabular}{llrrrr}
\toprule
landscape & condition & guessing, best k & guessing, three minus best & informed, best k & informed, three minus best \\
\midrule
grid & batch 5\% & 1 (24\%) & $-$0.19 [$-$0.22, $-$0.15] & 1 (23\%) & $-$0.17 [$-$0.21, $-$0.13] \\
grid & batch 10\% & 1 (24\%) & $-$0.15 [$-$0.19, $-$0.11] & 1 (23\%) & $-$0.14 [$-$0.18, $-$0.11] \\
grid & floor & 1 (24\%) & $-$0.17 [$-$0.21, $-$0.14] & 1 (23\%) & $-$0.17 [$-$0.20, $-$0.14] \\
grid & failed wells & 1 (24\%) & $-$0.22 [$-$0.26, $-$0.17] & 1 (23\%) & $-$0.21 [$-$0.25, $-$0.17] \\
\midrule
GB1 & batch 5\% & 0 (1\%) & $-$2.21 [$-$2.34, $-$2.08] & 0 (0\%) & $-$1.58 [$-$1.70, $-$1.47] \\
GB1 & batch 10\% & 0 (1\%) & $-$2.16 [$-$2.30, $-$2.03] & 0 (0\%) & $-$1.48 [$-$1.59, $-$1.36] \\
GB1 & floor & 0 (1\%) & $-$2.27 [$-$2.41, $-$2.14] & 0 (0\%) & $-$1.68 [$-$1.81, $-$1.55] \\
GB1 & failed wells & 0 (1\%) & $-$2.32 [$-$2.46, $-$2.16] & 0 (0\%) & $-$1.71 [$-$1.82, $-$1.59] \\
\bottomrule
\end{tabular}
\end{table}

